\documentclass[10pt, aspectratio=169]{beamer}
% Preámbulo modular para presentaciones Beamer (Modelación Estadística)
% Diseño y paleta institucional del Tecnológico de Monterrey

\documentclass[aspectratio=169,xcolor={dvipsnames,table}]{beamer}

\usepackage[utf8]{inputenc}
\usepackage[spanish,mexico]{babel}
\usepackage{amsmath,amssymb,amsfonts,mathtools}
\usepackage{graphicx}
\usepackage{listings}
\usepackage{booktabs}
\usepackage{tikz}
\usetikzlibrary{matrix,arrows,decorations.pathmorphing,shapes.geometric,positioning}

% Paleta Institucional Tecnológico de Monterrey
\definecolor{TecAzul}{HTML}{0039A6}       % Azul TEC: color primario institucional
\definecolor{TecAzulOscuro}{HTML}{1A2E51} % Azul marino oscuro
\definecolor{TecRojo}{HTML}{EC2661}       % Rojo de acento (paleta miTec)
\definecolor{TecGris}{HTML}{646464}       % Gris neutro para comentarios y subtítulos
\definecolor{TecFondoBloque}{HTML}{F4F6F9}% Fondo suave para bloques

% Configuración de Tema Beamer
\usetheme{Boadilla}
\usecolortheme[named=TecAzulOscuro]{structure}
\setbeamercolor{alerted text}{fg=TecRojo}
\setbeamercolor{palette primary}{bg=TecAzulOscuro,fg=white}
\setbeamercolor{palette secondary}{bg=TecAzul,fg=white}
\setbeamercolor{palette tertiary}{bg=TecAzulOscuro!80!black,fg=white}
\setbeamercolor{title}{fg=TecAzulOscuro,bg=TecFondoBloque}
\setbeamercolor{frametitle}{fg=TecAzulOscuro,bg=TecFondoBloque}

% Bloques personalizados elegantes
\setbeamercolor{block title}{bg=TecAzulOscuro,fg=white}
\setbeamercolor{block body}{bg=TecFondoBloque,fg=black}
\setbeamercolor{block title alerted}{bg=TecRojo,fg=white}
\setbeamercolor{block body alerted}{bg=TecRojo!8,fg=black}
\setbeamercolor{block title example}{bg=TecAzul,fg=white}
\setbeamercolor{block body example}{bg=TecAzul!8,fg=black}

% Redondear bordes y añadir sombra sutil
\setbeamertemplate{blocks}[rounded][shadow=true]
\setbeamertemplate{navigation symbols}{}

% Pie de página limpio con número de diapositiva
\setbeamertemplate{footline}{
  \leavevmode%
  \hbox{%
  \begin{beamercolorbox}[wd=.7\paperwidth,ht=2.25ex,dp=1ex,left]{author in head/foot}%
    \hspace*{2ex}\insertshortauthor~~(\insertshortinstitute) \hfill \insertshorttitle
  \end{beamercolorbox}%
  \begin{beamercolorbox}[wd=.3\paperwidth,ht=2.25ex,dp=1ex,right]{date in head/foot}%
    \insertshortdate{}\hspace*{2em}
    \insertframenumber{} / \inserttotalframenumber\hspace*{2ex}
  \end{beamercolorbox}}%
  \vskip0pt%
}

% Configuración de Listings de Python para visualización pedagógica de código
\lstset{
    language=Python,
    basicstyle=\ttfamily\footnotesize,
    keywordstyle=\color{TecAzulOscuro}\bfseries,
    stringstyle=\color{TecRojo},
    commentstyle=\color{TecGris}\itshape,
    showstringspaces=false,
    breaklines=true,
    frame=lines,
    rulecolor=\color{TecAzulOscuro!30},
    backgroundcolor=\color{TecFondoBloque},
    aboveskip=0.5em,
    belowskip=0.5em,
    literate={á}{{\'a}}1 {é}{{\'e}}1 {í}{{\'i}}1 {ó}{{\'o}}1 {ú}{{\'u}}1
             {Á}{{\'A}}1 {É}{{\'E}}1 {Í}{{\'I}}1 {Ó}{{\'O}}1 {Ú}{{\'U}}1
             {ñ}{{\~n}}1 {Ñ}{{\~N}}1 {¿}{{?`}}1 {¡}{{!`}}1
}

% Metadatos institucionales por defecto
\title[Modelación Estadística]{Modelación Estadística para Ciencia de Datos e Ingeniería}
\author[J. Castillo Colmenares]{Juliho David Castillo Colmenares}
\institute[Tec de Monterrey]{Tecnológico de Monterrey \\ Escuela de Ingeniería y Ciencias}
\date{\today}

% Comandos y macros estadísticas compartidas para las presentaciones Beamer (Metropolis)

% Conjuntos numéricos básicos
\newcommand{\R}{\mathbb{R}}
\newcommand{\N}{\mathbb{N}}
\newcommand{\Q}{\mathbb{Q}}
\newcommand{\Z}{\mathbb{Z}}
\newcommand{\set}[1]{\left\{ #1 \right\}}
\newcommand{\abs}[1]{\left|#1\right|}
\newcommand{\norm}[1]{\left\| #1 \right\|}

% Comandos estadísticos y probabilísticos del curso
\newcommand{\Var}{\operatorname{Var}}
\newcommand{\cov}{\operatorname{Cov}}
\newcommand{\corr}{\rho}
\newcommand{\comb}[2]{\begin{pmatrix} #1 \\ #2 \end{pmatrix}}
\newcommand{\s}{\sigma}
\newcommand{\particion}{\mathfrak{P}}
\newcommand{\card}[1]{\#\left( #1 \right)}
\newcommand{\seq}[3]{\set{#1}_{#2}^{#3}}
\newcommand{\E}{\mathbb{E}}
\newcommand{\Prob}{\mathbb{P}}

% Entornos pedagógicos y cajas en español académico natural
\newenvironment{discusion}{%
  \begin{alertblock}{Pregunta de Discusión}%
}{%
  \end{alertblock}%
}

\newenvironment{reflexion}{%
  \begin{alertblock}{Reflexión para el Aula}%
}{%
  \end{alertblock}%
}

\newenvironment{paradoja}{%
  \begin{alertblock}{Paradoja de Exploración}%
}{%
  \end{alertblock}%
}

\newenvironment{motivacion}{%
  \begin{exampleblock}{Motivación: Ciencia de Datos en la Práctica}%
}{%
  \end{exampleblock}%
}

\newenvironment{retoclase}{%
  \begin{block}{Reto Rápido en Equipo}%
}{%
  \end{block}%
}


\title[Distribución de Poisson]{Sección 03.07: Distribución de Poisson}
\subtitle{Procesos de Llegada, Ley de Eventos Raros y Aditividad}
\author[J. Castillo Colmenares]{Juliho Castillo Colmenares}
\institute{Tecnológico de Monterrey}
\date{\vspace{-1.2cm}}

\begin{document}

% Slide 1: Portada
\begin{frame}[plain]
  \vspace{-0.8cm}
  \titlepage
\end{frame}

% Slide 2: Hoja de Ruta Modular
\begin{frame}{Hoja de Ruta --- Capítulo 03: Variables Aleatorias Discretas}
  \footnotesize
  ¿Dónde nos encontramos en el desarrollo modular de la teoría? \pause
  \begin{itemize}\itemsep=0.08em
    \item \textbf{03.01} Concepto, Notación y Definición de Variable Aleatoria \pause
    \item \textbf{03.02} Funciones de Probabilidad (PMF) y Distribución Acumulada (CDF) \pause
    \item \textbf{03.03} Esperanza Matemática y Varianza Operativa \pause
    \item \textbf{03.04} Distribución Binomial y Ensayos de Bernoulli \pause
    \item \textbf{03.05} Distribuciones Geométrica y Binomial Negativa \pause
    \item \textbf{03.06} Distribución Hipergeométrica \pause
    \item \color{TecRojo}\textbf{03.07 Distribución de Poisson \emph{(Hoy)}}\color{black} \pause
    \item \textbf{03.08} Distribución Multinomial
  \end{itemize}
  \vspace{-0.1cm}
  \begin{block}{Objetivo de la Sesión}
    Modelar el conteo de eventos raros en intervalos de tiempo o espacio, dominar la ley de eventos raros como puente entre Binomial y Poisson, y explotar las propiedades de aditividad y la distribución condicional binomial bajo marginales fijos.
  \end{block}
\end{frame}

% Slide 3: Motivación
\begin{frame}{De la Binomial Discreta al Proceso Continuo de Poisson}
  \footnotesize
  \begin{motivacion}[¿Por qué una nueva distribución si ya tenemos la Binomial?]
    La distribución Binomial modela un número fijo de ensayos $n$ con probabilidad constante $p$. Pero en problemas reales (llamadas a un servidor, defectos en manufactura, mutaciones genéticas), los eventos ocurren de manera continua y aleatoria a una tasa $\lambda$ fija, sin que $n$ esté predefinido. El modelado correcto es entonces contar eventos en un intervalo, no ensayos en una muestra.
  \end{motivacion}

  \vspace{0.15cm}
  \pause
  \begin{itemize}\itemsep=0.1cm
    \item \textbf{Conteo continuo:} $X$ cuenta el número de eventos en un intervalo fijo de tiempo o espacio, no el número de éxitos en $n$ ensayos. \pause
    \item \textbf{Tasa media $\lambda$:} La esperanza y la varianza coinciden en una sola cantidad: $\E[X] = \Var(X) = \lambda$ (equidispersión). \pause
    \item \textbf{Soporte infinito numerable:} $\mathcal{S}_X = \set{0, 1, 2, \dots}$, adecuado para conteos sin cota superior fija.
  \end{itemize}
\end{frame}

% Slide 4: Definición Formal de la PMF
\begin{frame}{Definición Formal de la Distribución de Poisson}
  \footnotesize
  Una variable aleatoria discreta $X$ sigue una ley de Poisson con parámetro de tasa $\lambda > 0$ si su función de masa de probabilidad está dada por: \pause

  \vspace{0.1cm}
  \begin{block}{Función de Masa de Probabilidad (PMF)}
    $$P(X = k) = \dfrac{e^{-\lambda} \lambda^{k}}{k!}, \quad k \in \set{0, 1, 2, \dots}$$
  \end{block}

  \vspace{0.1cm}
  \pause
  \begin{alertblock}{Interpretación y Parametrización en Python}
    \begin{itemize}\itemsep=0.06cm
      \item \textbf{Factores:} El factor $e^{-\lambda}$ normaliza la suma; $\lambda^{k}/k!$ codifica la cola de la serie de Taylor.
      \item \textbf{Normalización:} $\sum_{k=0}^{\infty} \dfrac{e^{-\lambda} \lambda^{k}}{k!} = e^{-\lambda} \cdot e^{\lambda} = 1$.
      \item \textbf{En SciPy:} Se utiliza \texttt{poisson.pmf(k, mu=lambda)} y \texttt{poisson.cdf(k, mu=lambda)}.
    \end{itemize}
  \end{alertblock}
\end{frame}

% Slide 5: Equidispersión
\begin{frame}{Equidispersión: $\mu = \sigma^{2} = \lambda$}
  \footnotesize
  La distribución de Poisson tiene una propiedad estructural única: la varianza coincide exactamente con la esperanza, ambas iguales al parámetro de tasa $\lambda$. \pause

  \vspace{0.1cm}
  \begin{columns}[T]
    \begin{column}{0.48\textwidth}
      \begin{block}{Esperanza Matemática}
        $$\E[X] = \sum_{k=0}^{\infty} k \cdot \dfrac{e^{-\lambda} \lambda^{k}}{k!} = \lambda$$
        Se obtiene derivando la identidad $e^{\lambda} = \sum_{k=0}^{\infty} \lambda^{k}/k!$ respecto a $\lambda$ y multiplicando por $\lambda$.
      \end{block}
    \end{column}
    \begin{column}{0.48\textwidth}
      \begin{alertblock}{Varianza por Momento Factorial}
        $$\Var(X) = \E[X^{2}] - (\E[X])^{2} = (\lambda^{2} + \lambda) - \lambda^{2} = \lambda$$
        El momento factorial $\E[X(X-1)] = \lambda^{2}$ y la suma $\E[X^{2}] = \lambda^{2} + \lambda$ producen varianza idéntica a la media.
      \end{alertblock}
    \end{column}
  \end{columns}
\end{frame}

% Slide 6: Soporte y Normalización
\begin{frame}{Soporte Infinito Numerable y Normalización Exponencial}
  \footnotesize
  La suma de las probabilidades de Poisson sobre todo el soporte no es trivial: requiere la identidad de la serie de Taylor de la exponencial. \pause

  \vspace{0.1cm}
  \begin{block}{Soporte y Sumatoria Unitaria}
    $$\mathcal{S}_X = \set{k \in \mathbb{Z}_{\ge 0}}, \quad \sum_{k=0}^{\infty} P(X = k) = e^{-\lambda} \sum_{k=0}^{\infty} \dfrac{\lambda^{k}}{k!} = e^{-\lambda} \cdot e^{\lambda} = 1$$
  \end{block}

  \vspace{0.1cm}
  \pause
  \begin{alertblock}{Función Generadora de Probabilidades (PGF)}
    La PGF $G_{X}(t) = \E[t^{X}] = \sum_{k=0}^{\infty} t^{k} P(X=k) = e^{\lambda(t-1)}$ codifica toda la información distribucional y será clave para demostrar la aditividad.
  \end{alertblock}
\end{frame}

% Slide 7: Función Generadora de Momentos
\begin{frame}{MGF y Caracterización de la Distribución de Poisson}
  \footnotesize
  La Función Generadora de Momentos $M_{X}(t) = \E[e^{tX}]$ para una variable de Poisson adopta una forma cerrada especialmente simple: \pause

  \vspace{0.1cm}
  \begin{block}{MGF de Poisson}
    $$M_{X}(t) = \sum_{k=0}^{\infty} \dfrac{e^{tk} e^{-\lambda} \lambda^{k}}{k!} = e^{-\lambda} \sum_{k=0}^{\infty} \dfrac{(\lambda e^{t})^{k}}{k!} = e^{\lambda(e^{t} - 1)}$$
  \end{block}

  \vspace{0.1cm}
  \pause
  \begin{alertblock}{Derivación de Momentos desde la MGF}
    Los momentos brutos se obtienen por diferenciación: $\E[X^{r}] = M_{X}^{(r)}(0)$. Por ejemplo, $\E[X] = M_{X}'(0) = \lambda e^{0} = \lambda$ y $\E[X^{2}] = \lambda^{2} + \lambda$, confirmando la equidispersión.
  \end{alertblock}
\end{frame}

% Slide 8: Ley de Eventos Raros
\begin{frame}{Ley de Eventos Raros: Aproximación Binomial $\to$ Poisson}
  \footnotesize
  Cuando $n$ es muy grande y $p$ es muy pequeño, manteniendo el producto $np = \lambda$ constante, la distribución $\text{Bin}(n, p)$ converge puntualmente a $\text{Pois}(\lambda)$: \pause

  \vspace{0.1cm}
  \begin{block}{Límite de Poisson para Eventos Raros}
    $$\lim_{n \to \infty} \binom{n}{k} p^{k} (1-p)^{n-k} = \dfrac{e^{-\lambda} \lambda^{k}}{k!} \quad \text{cuando } np = \lambda \text{ constante}$$
  \end{block}

  \vspace{0.1cm}
  \pause
  \begin{alertblock}{Regla Operativa de Ingeniería}
    La aproximación es excelente bajo la regla práctica $n \ge 50$ y $p \le 0.05$ (con $np \le 5$). Por ejemplo, $n = 2000$ pacientes con $p = 0.001$ se aproximan por $\text{Pois}(2)$ con error absoluto $< 5 \times 10^{-4}$.
  \end{alertblock}
\end{frame}

% Slide 9: Aditividad
\begin{frame}{Aditividad: Suma de Poisson Independientes}
  \footnotesize
  Si $X \sim \text{Pois}(\lambda_{1})$ y $Y \sim \text{Pois}(\lambda_{2})$ son independientes, su suma $Z = X + Y$ también sigue una ley de Poisson con parámetro $\lambda_{1} + \lambda_{2}$. \pause

  \vspace{0.1cm}
  \begin{columns}[T]
    \begin{column}{0.48\textwidth}
      \begin{block}{Demostración por Convolución}
        \scriptsize
        $$P(Z = m) = \sum_{k=0}^{m} P(X=k) P(Y=m-k)$$
        $$= \dfrac{e^{-(\lambda_{1}+\lambda_{2})}(\lambda_{1}+\lambda_{2})^{m}}{m!}$$
        Aplicando el Teorema del Binomio de Newton a la suma interna de la convolución.
      \end{block}
    \end{column}
    \begin{column}{0.48\textwidth}
      \begin{alertblock}{Propiedad Reproductiva}
        \scriptsize
        La suma de $n$ variables Poisson independientes con tasas $\lambda_{1}, \dots, \lambda_{n}$ es $\text{Pois}(\sum \lambda_{i})$. Esta propiedad permite agregar flujos de eventos (e.g., dos servidores) sin perder la estructura.
      \end{alertblock}
    \end{column}
  \end{columns}
\end{frame}

% Slide 10: Propiedad Condicional Binomial
\begin{frame}{Propiedad Condicional: Distribución Binomial bajo Suma Fija}
  \footnotesize
  Si $X \sim \text{Pois}(\lambda_{1})$ y $Y \sim \text{Pois}(\lambda_{2})$ son independientes, entonces la distribución de $X$ condicionada al total $X + Y = n$ es una Binomial: \pause

  \vspace{0.1cm}
  \begin{block}{Distribución Condicional Binomial}
    $$P(X = k \mid X + Y = n) = \binom{n}{k} \left(\dfrac{\lambda_{1}}{\lambda_{1} + \lambda_{2}}\right)^{k} \left(\dfrac{\lambda_{2}}{\lambda_{1} + \lambda_{2}}\right)^{n-k}, \quad k = 0, 1, \dots, n$$
  \end{block}

  \vspace{0.1cm}
  \pause
  \begin{alertblock}{Interpretación en Teoría de Colas}
    Si dos servidores atienden clientes Poisson con tasas $\lambda_{1}$ y $\lambda_{2}$, entonces dado que se atendieron en total $n$ clientes, el número que atendió el servidor $1$ sigue una $\text{Bin}(n, \lambda_{1}/(\lambda_{1}+\lambda_{2}))$, donde $p$ es la fracción de tráfico asignada.
  \end{alertblock}
\end{frame}

% Slide 11: Aplicaciones
\begin{frame}{Aplicaciones en Industria y Ciencia de Datos}
  \footnotesize
  El modelado de Poisson es la piedra angular de tres pilares modernos de la ingeniería y la analítica de datos: \pause

  \vspace{0.15cm}
  \begin{itemize}\itemsep=0.15cm
    \item \textbf{Teoría de Colas y Telecomunicaciones:} Modelado de llegadas de llamadas a un call center, peticiones HTTP a un servidor, o vehículos a una caseta de peaje. La tasa $\lambda$ define la intensidad de tráfico. \pause
    \item \textbf{Control de Calidad Industrial:} Conteo de defectos por unidad de producción, donde la equidispersión permite cartas de control $c$ con límites $\lambda \pm 3\sqrt{\lambda}$. \pause
    \item \textbf{Bioestadística y Epidemiología:} Número de mutaciones genéticas en una cadena de ADN, casos raros de una enfermedad en una población, o eventos adversos raros en ensayos clínicos.
  \end{itemize}
\end{frame}

% Slide 12A-1: Lab Python (Bloque 1)
\begin{frame}[fragile]{Laboratorio en Python: PMF y Equidispersión}
  \scriptsize
  Verificación computacional en SciPy (\texttt{03.07\_poisson\_distribution.py}) del soporte, normalización y propiedad de equidispersión:
  \vspace{0.08cm}
  {\lstset{basicstyle=\fontsize{5.5pt}{6.5pt}\selectfont\ttfamily}
  \lstinputlisting[firstline=15, lastline=37]{../../code/03_variables_aleatorias_discretas/03.07_poisson_distribution.py}}
\end{frame}

% Slide 12A-2: Lab Python (Bloque 2)
\begin{frame}[fragile]{Laboratorio en Python: Ley de Eventos Raros}
  \scriptsize
  Simulación de $250,000$ ensayos verificando la convergencia Binomial $\to$ Poisson mediante Distancia de Variación Total (TVD):
  \vspace{0.05cm}
  {\lstset{basicstyle=\fontsize{5.5pt}{6.5pt}\selectfont\ttfamily}
  \lstinputlisting[firstline=40, lastline=62]{../../code/03_variables_aleatorias_discretas/03.07_poisson_distribution.py}}
\end{frame}

% Slide 12B: Lab Python (Bloque 3)
\begin{frame}[fragile]{Laboratorio en Python: Aditividad y Condicional Binomial}
  \scriptsize
  Verificación numérica por simulación de la suma de Poisson y de la distribución condicional binomial:
  \vspace{0.02cm}
  {\lstset{basicstyle=\fontsize{5pt}{6pt}\selectfont\ttfamily}
  \lstinputlisting[firstline=72, lastline=92]{../../code/03_variables_aleatorias_discretas/03.07_poisson_distribution.py}}
\end{frame}

% Slide 12C: Salida Terminal
\begin{frame}[fragile]{Laboratorio en Python: Salida en Terminal}
  \vspace{-0.1cm}
  \begin{block}{}
    \fontsize{4.5pt}{5.2pt}\selectfont
    \begin{verbatim}
=== Block 1: Poisson PMF & Equidispersion Validation ===
Lambda = 3.0 (calls/min) | mu = sigma^2 = lambda
PMF total probability: 0.999999329614
Equidispersion property verified: True
P(X=0)=0.049787 | P(X=1)=0.149361 | P(X=2)=0.224042

=== Block 2: Law of Rare Events (Binomial -> Poisson) ===
  n=    50 p=0.0400 | TVD=0.009270
  n=   500 p=0.0040 | TVD=0.000905
  n=  2000 p=0.0010 | TVD=0.000226
  n= 10000 p=0.0002 | TVD=0.000045
Monte Carlo (N=250,000): mean=2.0018 | var=1.9971

=== Block 3: Additivity & Conditional Binomial ===
Z = X + Y: emp mean=7.9992 | emp var=7.9888
Conditional X | (X+Y=8) with p_cond = 0.3750:
  k=0: 0.0222 | 0.0233 | k=3: 0.2856 | 0.2816
  k=5: 0.1019 | 0.1014 | k=8: 0.0003 | 0.0004
    \end{verbatim}
  \end{block}
\end{frame}

% Slide 17: Cierre
\begin{frame}{Cierre de Sección y Síntesis sobre la Ley de Poisson}
  \begin{reflexion}[El Poder de los Eventos Raros en la Ingeniería]
    La distribución de Poisson nos enseña que, en sistemas donde los eventos son independientes y ocurren a tasa constante, tanto la incertidumbre como el valor esperado se cuantifican con un solo parámetro $\lambda$. Esta propiedad de equidispersión es la base del control estadístico de procesos y del modelado de colas.
  \end{reflexion}

  \vspace{0.15cm}
  \pause
  \begin{block}{Lecciones Clave de la Ley de Eventos Raros}
    \begin{itemize}\itemsep=0.08cm
      \item \textbf{Equidispersión:} $\E[X] = \Var(X) = \lambda$ permite usar la misma magnitud para predicción y monitoreo de variabilidad.
      \item \textbf{Aditividad:} La suma de Poisson independientes es Poisson; ideal para agregar flujos de eventos en redes de servicio.
      \item \textbf{Condicional Binomial:} Bajo suma fija, la distribución de una componente es binomial; permite asignación proporcional de tráfico entre servidores.
    \end{itemize}
  \end{block}
\end{frame}

% Slide 18: Perspectiva Modular
\begin{frame}{Perspectiva Modular --- El Próximo Reto}
  \scriptsize
  \begin{retoclase}[Para la próxima sesión: Distribución Multinomial]
    Hasta ahora hemos estudiado distribuciones para variables aleatorias que cuentan un solo tipo de evento (éxitos, defectos, llamadas, clientes). Pero, ¿qué sucede cuando un ensayo puede resultar en $k \ge 2$ categorías mutuamente excluyentes con probabilidades $p_{1}, p_{2}, \dots, p_{k}$ que suman $1$? Introduciremos la \emph{Distribución Multinomial}.
  \end{retoclase}

  \vspace{0.08cm}
  \pause
  \begin{alertblock}{Preparación Recomendada}
    \begin{itemize}\itemsep=0.06cm
      \item Repasa la estructura del coeficiente multinomial $\binom{n}{x_{1}, \dots, x_{k}} = \dfrac{n!}{x_{1}! \cdots x_{k}!}$.
      \item Reflexiona sobre cómo modelarías las probabilidades conjuntas de las preferencias de voto de una población entre $k$ candidatos.
    \end{itemize}
  \end{alertblock}
\end{frame}

\end{document}
